\documentclass[conference]{IEEEtran}

% ---------------------------------------------------------------
% Packages
% ---------------------------------------------------------------
\usepackage{graphicx}
\usepackage{booktabs}
\usepackage{amsmath}
\usepackage{amssymb}
\usepackage{cite}
\usepackage{url}
\usepackage{array}
\usepackage{multirow}

% ---------------------------------------------------------------
% Title and Authors
% ---------------------------------------------------------------
\title{Trainable Nonlinear Reaction Diffusion for Multiplicative Gamma Noise Removal:\\
A Stage-Wise Learned Model with PDE Baseline Comparison}

\author{%
  \IEEEauthorblockN{Mahipal Jetta}
  \IEEEauthorblockA{Mahindra University}
  \and
  \IEEEauthorblockN{Sujato Dutta}
  \IEEEauthorblockA{Mahindra University}
}

% ---------------------------------------------------------------
\begin{document}
\maketitle

% ---------------------------------------------------------------
\begin{abstract}
This paper presents an extensive study on the application of Trainable Nonlinear
Reaction Diffusion (TNRD) for denoising grayscale images corrupted by multiplicative
gamma noise. The proposed model consists of a sequential reaction--diffusion update
in which the filter kernels are fixed (randomly initialized), while the nonlinear
influence functions and regularization coefficients are learnable. The model is
compared against a nonlinear smooth diffusion PDE specifically designed for gamma
noise. Experiments are conducted on a 400-image grayscale dataset under two noise
levels, $L=1$ and $L=10$. Performance is measured using PSNR and SSIM, complemented
by qualitative analysis and ablation studies. Results show that the proposed TNRD
model achieves the best mean PSNR for the more challenging case of $L=1$,
outperforming the PDE baseline by approximately $0.49$~dB, whereas the PDE model
performs slightly better in both mean PSNR and SSIM for the milder case of $L=10$.
Ablation experiments confirm that each component --- the nonlinear influence
functions, the full filter bank, and the multi-stage design --- contributes
meaningfully to overall performance.
\end{abstract}

\begin{IEEEkeywords}
image denoising, multiplicative gamma noise, trainable nonlinear reaction diffusion,
PDE baseline, grayscale image restoration, stage-wise learning
\end{IEEEkeywords}

% ---------------------------------------------------------------
\section{Introduction}
% ---------------------------------------------------------------
Multiplicative gamma noise is an important class of image degradation in which
the observed intensity depends on the local signal magnitude. Unlike additive
Gaussian noise, multiplicative noise affects local contrast in proportion to signal
strength, making it particularly challenging to handle. Conventional smoothing
methods tend to over-smooth low-intensity regions or lose high-frequency detail,
while classical diffusion-based techniques typically rely on finely tuned,
hand-crafted coefficients.

The Trainable Nonlinear Reaction--Diffusion (TNRD) framework \cite{chen2017tnrd}
offers a principled middle ground between purely model-based and purely
learning-based restoration. Rather than learning a large, general-purpose black
box, TNRD models the restoration process as a sequence of diffusion steps, each
capable of learning its own nonlinear dynamics, while retaining the interpretable
structure of a variational/PDE-based formulation.

The central contribution of this paper is to apply TNRD to multiplicative gamma
noise and to compare it directly against a nonlinear smooth diffusion PDE baseline
designed for the same noise model. The goals are twofold: to assess whether the
learned model surpasses the hand-crafted PDE, and to isolate the contribution of
each architectural component through ablation studies.

The specific contributions of this paper are as follows.
\begin{itemize}
  \item We train a multi-stage TNRD denoiser for multiplicative gamma noise
        reduction using a structured update equation with fixed random filters and
        per-stage learned nonlinear influence functions and scalar regularization
        weights.
  \item We adopt a nonlinear smooth diffusion PDE, guided by both image intensity
        and gradient magnitude, as a strong and principled baseline that improves
        upon general isotropic models and the standard Perona--Malik approach.
  \item Experiments are performed under two gamma noise regimes ($L=1$ and $L=10$)
        on a 400-image grayscale dataset.
  \item Ablation studies quantify the individual contribution of the nonlinear
        influence functions, filter bank size, and number of stages.
  \item A full reproducible pipeline --- including checkpoints, evaluation metrics,
        figures, and an IEEE-formatted report --- is provided.
\end{itemize}

% ---------------------------------------------------------------
\section{Related Work}
% ---------------------------------------------------------------
Diffusion-based image restoration has a long history. Perona and Malik introduced
anisotropic diffusion, demonstrating that the smoothing process can be made
edge-aware by varying the diffusion coefficient according to local gradient
magnitude \cite{perona1990scale}. Chen and Pock later proposed the TNRD framework,
which forged an explicit link between variational restoration and learned
reaction--diffusion models, and achieved strong restoration performance with high
interpretability \cite{chen2017tnrd}.

For multiplicative noise, PDE-based approaches are especially attractive because
their spatially adaptive diffusion coefficients can be conditioned on both image
intensity and gradient magnitude --- a natural fit for gamma-noise and speckle-like
corruption, where diffusion should behave differently in uniform regions, near
edges, and in low-contrast areas. The nonlinear smooth diffusion PDE adopted as the
baseline in this work exploits precisely these properties.

In terms of novelty, this paper does not propose a new theoretical diffusion model.
Instead, it applies the TNRD framework to multiplicative gamma noise under a
specific constraint: the convolutional filters are kept fixed while the nonlinear
influence functions and regularization coefficients are estimated in a stage-wise
fashion.

% ---------------------------------------------------------------
\section{Problem Formulation}
% ---------------------------------------------------------------
Let $u$ denote the clean grayscale image and $f$ the noisy observation. The
multiplicative gamma noise model is
\begin{equation}
  f = u \cdot n, \qquad n \sim \Gamma(L,\,L),
  \label{eq:noise_model}
\end{equation}
where $L$ is the shape parameter of the gamma distribution. Two settings are
evaluated: $L=1$ (severe noise) and $L=10$ (moderate noise).

The proposed TNRD model performs the stage-wise update
\begin{align}
  u_t = u_{t-1} - \Bigg(
    &\sum_{i=1}^{N_k} \bar{k}_i^{\,t} * \!\left(
        \frac{u_\sigma}{M}\,\phi_i^t\!\left(k_i^t * u_{(t-1)p}\right)
      \right) \nonumber \\
    &+ \lambda\,\frac{u - f}{u^2 + \varepsilon}
  \Bigg),
  \label{eq:tnrd_update}
\end{align}
where:
\begin{itemize}
  \item $u_t$ is the restored image at stage $t$;
  \item $u_{(t-1)p}$ is the zero-padded output of the previous stage;
  \item $u_\sigma$ is a Gaussian-smoothed version of the current estimate;
  \item $k_i^t$ are fixed convolution filters and $\bar{k}_i^{\,t}$ their
        spatially flipped counterparts used in the reverse diffusion step;
  \item $\phi_i^t$ are learnable influence functions;
  \item $\lambda$ is a learnable scalar regularization parameter; and
  \item $M$ is the gamma noise parameter used during training and evaluation.
\end{itemize}

In summary, the fixed filter bank extracts linear features, the learned influence
functions introduce stage-specific nonlinearity, and the reaction term modulates
data fidelity. Because only the influence functions and $\lambda$ are optimized,
the update remains mathematically interpretable.

% ---------------------------------------------------------------
\section{Proposed Method}
% ---------------------------------------------------------------

\subsection{TNRD Architecture}
The proposed architecture uses $N_k = 8$ filters of size $3\times3$ across five
stages. Each stage has its own fixed filter set and its own influence function.
Filter weights are sampled randomly at initialization and remain frozen throughout
training, in accordance with the design requirement that $k_i$ be fixed while all
other parameters are learned per stage.

Each influence function is implemented as a grouped $1\times1$ convolution followed
by a hyperbolic tangent activation, allowing each filter response to be transformed
individually through a trainable nonlinear mapping. The scalar $\lambda$ is also
learned independently per stage, enabling the diffusion--fidelity trade-off to
evolve across the unrolled iterations.

Gaussian blurring is applied with zero-padding to obtain $u_\sigma$. The forward
diffusion pass uses \texttt{F.conv2d}, while the reverse pass uses
\texttt{F.conv\_transpose2d} with spatially reversed filters.

\subsection{Stage-Wise Learning Strategy}
The network is trained sequentially: once a stage has been optimized, its parameters
are frozen and only the parameters of the next stage are updated. This greedy
layer-wise scheme, analogous to the original TNRD training protocol, avoids
simultaneous optimization across all stages and improves training stability.

The loss function is the mean squared error (MSE) between the restored image and
the clean reference. Training uses the Adam optimizer with a learning rate of
$10^{-3}$, a batch size of $1$, and $30$ epochs for the full network.

\subsection{Novel Contribution}
The principal novelty of this work is the application of TNRD to a multiplicative
gamma noise model under a fixed-filter constraint, with stage-wise nonlinearity
learning. Unlike conventional deep denoising networks, the proposed approach is
grounded in an explicit mathematical update rule that makes the interaction between
diffusion and reaction interpretable. This is further supported by a direct
comparison with a nonlinear diffusion PDE designed for the same noise model.

% ---------------------------------------------------------------
\section{PDE Baseline Model}
% ---------------------------------------------------------------
The baseline is a nonlinear smooth diffusion PDE of the form
\begin{equation}
  \frac{\partial u}{\partial t}
    = \operatorname{div}\!\left(g(u_\sigma,\,|\nabla u_\sigma|)\,\nabla u\right),
  \label{eq:pde}
\end{equation}
with diffusion coefficient
\begin{equation}
  g(u_\sigma,\,|\nabla u_\sigma|)
    = \left(\frac{u_\sigma}{M}\right)^{\!\alpha}
      \frac{1}{1 + |\nabla u_\sigma|^{\,\beta}}.
  \label{eq:diffusion_coeff}
\end{equation}
The factor $(u_\sigma/M)^\alpha$ scales diffusion with local intensity, reducing
smoothing in bright regions where multiplicative noise is strongest. The denominator
suppresses diffusion near edges, where large gradient magnitudes are present.
Numerically, \eqref{eq:pde} is solved via Gaussian pre-smoothing, directional
finite differences, face-centred averaging of diffusion coefficients, and explicit
time integration with replicated boundary conditions.

This baseline constitutes a strong benchmark because it is purpose-built for
multiplicative gamma noise using domain knowledge, so any improvement over it
reflects the genuine benefit of the learned TNRD components.

% ---------------------------------------------------------------
\section{Experimental Setup}
% ---------------------------------------------------------------

\subsection{Dataset and Split}
Experiments use a grayscale image dataset of 400 images from the
\texttt{FoETrainingSets180} directory. Images are discovered by recursive
traversal of the data folder and partitioned into 80\% training, 10\% validation,
and 10\% test sets, yielding 40 test images.

\subsection{Noise Settings}
Two multiplicative gamma noise configurations are evaluated:
\begin{itemize}
  \item Severe noise: $L = 1$
  \item Moderate noise: $L = 10$
\end{itemize}
These settings allow performance to be assessed across a wide range of noise
intensities.

\subsection{Evaluation Metrics}
Performance is measured using Peak Signal-to-Noise Ratio (PSNR, in dB) and the
Structural Similarity Index (SSIM) \cite{wang2004ssim}. PSNR captures pixel-level
distortion, while SSIM reflects perceptual and structural fidelity. Both metrics
are computed for the noisy input, the PDE output, and the proposed TNRD output.

\subsection{Ablation Variants}
Four TNRD configurations are evaluated to isolate the contribution of each design
choice:
\begin{itemize}
  \item \textbf{Full model}: $N_k = 8$ filters, 5 stages, nonlinear influence
        functions.
  \item \textbf{No nonlinear influence}: $N_k = 8$ filters, 5 stages, identity
        influence functions (linear ablation).
  \item \textbf{Reduced filters}: $N_k = 4$ filters, 5 stages.
  \item \textbf{Reduced stages}: $N_k = 8$ filters, 3 stages.
\end{itemize}

% ---------------------------------------------------------------
\section{Results and Discussion}
% ---------------------------------------------------------------

\subsection{Quantitative Comparison with the PDE Baseline}
Table~\ref{tab:main_results} reports mean PSNR and SSIM over 40 test images for
both noise levels. The PDE baseline substantially outperforms the noisy input in
both settings, confirming its effectiveness. The proposed TNRD model achieves the
highest mean PSNR for $L=1$.

\begin{table}[t]
  \centering
  \caption{Average test performance over 40 images.}
  \label{tab:main_results}
  \begin{tabular}{llcc}
    \toprule
    Noise Level & Method & PSNR (dB) & SSIM \\
    \midrule
    \multirow{3}{*}{$L=1$}
      & Noisy          & 10.7128 & 0.1352 \\
      & PDE baseline   & 18.3465 & \textbf{0.4213} \\
      & Proposed TNRD  & \textbf{18.8389} & 0.3433 \\
    \midrule
    \multirow{3}{*}{$L=10$}
      & Noisy          & 18.1927 & 0.4304 \\
      & PDE baseline   & \textbf{23.7250} & \textbf{0.6595} \\
      & Proposed TNRD  & 23.4794 & 0.5889 \\
    \bottomrule
  \end{tabular}
\end{table}

For $L=1$, the proposed TNRD model exceeds the PDE baseline by approximately
$0.49$~dB in mean PSNR, demonstrating the ability of the learned nonlinear
reaction--diffusion process to handle severe multiplicative noise. The PDE baseline
nonetheless achieves a higher mean SSIM, indicating that the hand-crafted smoothing
approach better preserves global structural coherence.

For $L=10$, the PDE baseline retains a marginal mean PSNR advantage of approximately
$0.25$~dB and a more noticeable SSIM advantage, suggesting that in mild-noise
conditions the manually designed diffusion scheme is highly competitive and the
learned TNRD framework has not extracted sufficient additional information to
surpass it.

In both settings the proposed TNRD model achieves a higher per-image PSNR than the
PDE baseline on 21 out of 40 test images, indicating that TNRD is not uniformly
weaker but performs well on a specific subset of images consistent with the local,
nonlinear nature of its stage-wise updates.

\subsection{Qualitative Comparison}
Figures~\ref{fig:l1} and \ref{fig:l10} show representative results for each noise
level. At $L=1$, the PDE baseline removes a large proportion of noise but tends
to produce a slightly more diffused appearance, whereas the proposed TNRD model
often restores stronger local contrast. At $L=10$, both methods produce visually
plausible outputs with more subtle differences; the PDE output is generally
smoother while the TNRD output occasionally preserves finer texture.

\begin{figure}[t]
  \centering
  \includegraphics[width=\linewidth]{results/figures/comparison_L1.png}
  \caption{Qualitative comparison for severe gamma noise ($L=1$). From left to
    right: clean image, noisy image, nonlinear smooth diffusion PDE baseline, and
    proposed TNRD output.}
  \label{fig:l1}
\end{figure}

\begin{figure}[t]
  \centering
  \includegraphics[width=\linewidth]{results/figures/comparison_L10.png}
  \caption{Qualitative comparison for moderate gamma noise ($L=10$). Column order
    as in Fig.~\ref{fig:l1}.}
  \label{fig:l10}
\end{figure}

Figure~\ref{fig:psnr} summarizes the average PSNR comparison graphically. Two
trends are immediately visible: both restoration methods substantially outperform
the raw noisy input, and the relative advantage between PDE and TNRD reverses
depending on the noise regime.

\begin{figure}[t]
  \centering
  \includegraphics[width=\linewidth]{results/figures/psnr_comparison.png}
  \caption{Average PSNR for the noisy input, PDE baseline, and proposed TNRD
    output under both gamma noise settings.}
  \label{fig:psnr}
\end{figure}

% ---------------------------------------------------------------
\section{Ablation Study}
% ---------------------------------------------------------------
Table~\ref{tab:ablation} presents the ablation results. The full TNRD model is the
best-performing configuration under both noise levels, validating the combination
of the full filter bank, nonlinear influence functions, and five-stage depth.

\begin{table}[t]
  \centering
  \caption{Ablation results for the proposed TNRD model.}
  \label{tab:ablation}
  \begin{tabular}{lcccc}
    \toprule
    Variant & $L$ & Filters & Stages & PSNR (dB) \quad SSIM \\
    \midrule
    Full model              & 1  & 8 & 5 & \textbf{17.5626} \quad \textbf{0.2532} \\
    No nonlinear influence  & 1  & 8 & 5 & $\phantom{0}$7.4212 \quad 0.0129 \\
    $N_k = 4$               & 1  & 4 & 5 & 15.9145 \quad 0.2122 \\
    3 stages                & 1  & 8 & 3 & 15.7740 \quad 0.2174 \\
    \midrule
    Full model              & 10 & 8 & 5 & \textbf{22.9860} \quad \textbf{0.5713} \\
    No nonlinear influence  & 10 & 8 & 5 & 13.4083 \quad 0.1704 \\
    $N_k = 4$               & 10 & 4 & 5 & 19.3513 \quad 0.4462 \\
    3 stages                & 10 & 8 & 3 & 20.6764 \quad 0.4933 \\
    \bottomrule
  \end{tabular}
\end{table}

Removing the nonlinear influence functions causes the largest performance drop,
particularly for $L=1$, underscoring that learned nonlinearity is the most critical
component of the proposed TNRD model. Reducing filter bank size or stage count each
leads to a measurable degradation, confirming that both properties are necessary.
Overall, the ablation study demonstrates that no component of the proposed
architecture is redundant.

% ---------------------------------------------------------------
\section{Key Observations}
% ---------------------------------------------------------------
The following key observations emerge from this study.
\begin{itemize}
  \item The proposed TNRD model is competitive with the purpose-built PDE baseline,
        substantially outperforming the noisy input in both noise settings.
  \item TNRD shows its greatest advantage under severe noise ($L=1$), where its mean
        PSNR exceeds that of the PDE baseline.
  \item The PDE baseline is highly competitive and often superior in terms of mean
        SSIM, reflecting its superior preservation of global structural coherence.
  \item Nonlinear influence functions are critical; removing them causes a
        dramatic performance drop, confirming that learned nonlinearity is the key
        driver of TNRD's effectiveness.
  \item A model can simultaneously achieve higher PSNR and lower SSIM than a
        competing method, highlighting the complementary nature of these metrics.
\end{itemize}

% ---------------------------------------------------------------
\section{Limitations and Future Scope}
% ---------------------------------------------------------------
Several limitations warrant acknowledgment. First, the use of fixed random filters
limits the model's adaptability to the input data distribution. Second, the
proposed model modestly underperforms the PDE baseline under milder noise
($L=10$). Third, the study is restricted to grayscale images from a single dataset,
limiting generalizability.

Promising directions for future work include: (i) relaxing the fixed-filter
constraint to allow joint training of filters and influence functions; (ii)
exploring stronger parameterizations of the influence functions and increasing the
number of stages; and (iii) extending the evaluation to additional datasets,
alternative multiplicative noise models, and color image restoration.

% ---------------------------------------------------------------
\section{Conclusion}
% ---------------------------------------------------------------
This paper presents an extensive experimental evaluation of a TNRD model for
multiplicative gamma noise removal, compared against a nonlinear smooth diffusion
PDE baseline. The TNRD framework performs structured, interpretable updates using
fixed convolutional filters, learned nonlinear influence functions, and per-stage
scalar regularization parameters. Experiments show that TNRD achieves the best
mean PSNR in the high-noise setting ($L=1$), while the PDE model performs slightly
better in both mean PSNR and SSIM for the milder setting ($L=10$). Ablation studies
confirm that the full architecture is necessary, with the nonlinear influence
functions contributing most to overall performance. Together, these results
demonstrate that a structurally constrained yet data-driven TNRD model can compete
with a purpose-built PDE baseline, and establishes TNRD as a meaningful approach
for denoising images corrupted by severe multiplicative gamma noise.

% ---------------------------------------------------------------
\begin{thebibliography}{99}

\bibitem{perona1990scale}
P.~Perona and J.~Malik,
``Scale-space and edge detection using anisotropic diffusion,''
\emph{IEEE Trans.\ Pattern Anal.\ Mach.\ Intell.},
vol.~12, no.~7, pp.~629--639, Jul.\ 1990.

\bibitem{chen2017tnrd}
Y.~Chen and T.~Pock,
``Trainable nonlinear reaction diffusion: A flexible framework for fast and
effective image restoration,''
\emph{IEEE Trans.\ Pattern Anal.\ Mach.\ Intell.},
vol.~39, no.~6, pp.~1256--1272, Jun.\ 2017.

\bibitem{wang2004ssim}
Z.~Wang, A.~C.~Bovik, H.~R.~Sheikh, and E.~P.~Simoncelli,
``Image quality assessment: From error visibility to structural similarity,''
\emph{IEEE Trans.\ Image Process.},
vol.~13, no.~4, pp.~600--612, Apr.\ 2004.

\bibitem{goodman1976speckle}
J.~W.~Goodman,
``Some fundamental properties of speckle,''
\emph{J.\ Opt.\ Soc.\ Amer.},
vol.~66, no.~11, pp.~1145--1150, Nov.\ 1976.

\end{thebibliography}

\end{document}
